\section{Sibling-seed composition: off-diagonal cells from composed seeds}

RelMath's prose promises that predication cells with \(\text{subject} \neq
\text{object}\) ``arise only by composing such seeds drawn from sibling
utterances of the articulating recursion, never by postulating an entity from
nothing.'' This section discharges that promise, and states honestly the one
piece of the promise the corpus never glossed: what ``sibling'' means.

\begin{remark}[The linear-sibling hypothesis --- a stated choice, not a theorem]\label{rem:sibling-hypothesis}\uses{def:articulating}
The corpus never defines which utterances of an \texttt{Articulating}
recursion count as one another's siblings. The reading formalized below is
the \emph{linear} one: for a recursion \(a\), its siblings are the utterances
of \texttt{a.utterances}, the flat list obtained by walking the
\texttt{pair}/\texttt{utter} spine and collecting every utterance carried
along the way, in order. This is the natural reading of a linear recursion
(each \texttt{pair} has one utterance and one further articulating, no
branching), but it is a choice: a structural notion of ``sibling'' (for
instance, utterances at the same depth across distinct recursions) is not
ruled out and is not addressed here.
\end{remark}

\begin{definition}[The utterances of a recursion]\label{thm:articulating-utterances}\lean{Obsmetrologia.Articulating.utterances}\leanok\uses{def:articulating, rem:sibling-hypothesis}
\texttt{utterances} walks the \texttt{pair}/\texttt{utter} spine of an
\texttt{Articulating} value and returns the list of every utterance carried
along the way, in order: \(\texttt{utter}\,u \mapsto [u]\) and
\(\texttt{pair}\,u\,a \mapsto u :: a.\mathrm{utterances}\). Under the
linear-sibling hypothesis of \cref{rem:sibling-hypothesis} this is the pool a
recursion's own siblings are drawn from.
\end{definition}

\begin{proposition}[The head is among its own utterances]\label{thm:head-mem-utterances}\lean{Obsmetrologia.Articulating.head_mem_utterances}\leanok\uses{thm:articulating-utterances, thm:floor}
The head utterance of a recursion is always a member of its own
\texttt{utterances} list --- what \cref{thm:seed-head-mem-cellsof} draws on
below is honestly drawn from \(a\)'s own siblings, not smuggled in from
outside.
\end{proposition}

\begin{proposition}[Composing sibling seeds]\label{thm:composeseeds}\lean{Obsmetrologia.PredicationCell.composeSeeds}\leanok\uses{thm:predication-seed}
The composite of the reflexive seeds of two utterances \(u, v\), taken
through the \emph{first} utterance's claim: \(\texttt{composeSeeds}\,u\,v =
(u.\mathrm{token}, u.\mathrm{claim}, v.\mathrm{token})\). Subject and
predicate come from \(u\); the object comes only from \(v\)'s token.
Composition is not symmetric in its two arguments: in general \(v\)'s
predicate never enters.
\end{proposition}

\begin{proposition}[Composing a seed with itself is the diagonal]\label{thm:composeseeds-self}\lean{Obsmetrologia.PredicationCell.composeSeeds_self}\leanok\uses{thm:composeseeds}
\(\texttt{composeSeeds}\,u\,u = \texttt{seed}\,u\). Nothing off-diagonal is
produced from a single utterance alone.
\end{proposition}

\begin{proposition}[The composite reads its endpoints off the two seeds]\label{thm:composeseeds-subj-obj}\lean{Obsmetrologia.PredicationCell.composeSeeds_subject_object}\leanok\uses{thm:composeseeds}
The composite's subject is \(\texttt{seed}\,u\)'s subject and its object is
\(\texttt{seed}\,v\)'s object. This is the sense in which an off-diagonal
cell is assembled from two seeds rather than posited fresh.
\end{proposition}

\subsection*{Projection recovery: a real, stated asymmetry}

\begin{theorem}[Subject-projection recovers the composite exactly]\label{thm:subjproj-composeseeds}\lean{Obsmetrologia.PredicationCell.subjProj_composeSeeds}\leanok\uses{thm:composeseeds, thm:predication-proj}
\(\texttt{subjProj}\,(\texttt{composeSeeds}\,u\,v) = \texttt{seed}\,u\),
exactly, predicate included --- because the composite's predicate already is
\(u\)'s claim.
\end{theorem}

\begin{theorem}[Object-projection does not, in general, recover the composite]\label{thm:objproj-composeseeds}\lean{Obsmetrologia.PredicationCell.objProj_composeSeeds}\leanok\uses{thm:composeseeds, thm:predication-proj}
\(\texttt{objProj}\,(\texttt{composeSeeds}\,u\,v) = (v.\mathrm{token},
u.\mathrm{claim}, v.\mathrm{token})\): collapsing the composite onto its
object side keeps \emph{\(u\)'s} claim, not \(v\)'s. The true statement is
this asymmetric form. A symmetric claim --- that object-projection recovers
\(\texttt{seed}\,v\) outright, the mirror of
\cref{thm:subjproj-composeseeds} --- is \textbf{not} claimed and is false in
general (see \cref{thm:objproj-composeseeds-of-claim-eq} for the one
condition under which it does hold).
\end{theorem}

\begin{proposition}[What object-projection does share with \texttt{seed v}]\label{thm:objproj-composeseeds-endpoints}\lean{Obsmetrologia.PredicationCell.objProj_composeSeeds_endpoints}\leanok\uses{thm:objproj-composeseeds}
Object-projection of the composite agrees with \(\texttt{seed}\,v\) at both
entity endpoints (subject and object, both \(v.\mathrm{token}\)) --- but not,
by \cref{thm:objproj-composeseeds}, at the predicate coordinate. Stated at the
endpoint level rather than as full equality.
\end{proposition}

\begin{proposition}[The one condition recovering \texttt{seed v} exactly]\label{thm:objproj-composeseeds-of-claim-eq}\lean{Obsmetrologia.PredicationCell.objProj_composeSeeds_of_claim_eq}\leanok\uses{thm:objproj-composeseeds}
When \(u\) and \(v\) share a claim, object-projection of the composite does
equal \(\texttt{seed}\,v\) exactly. This is the missing case that closes the
gap \cref{thm:objproj-composeseeds} leaves open, not a retraction of the
asymmetry: outside this hypothesis the equality fails.
\end{proposition}

\subsection*{Completeness: nothing off-diagonal need be postulated}

\begin{theorem}[\texttt{composeSeeds} is surjective onto all predication cells]\label{thm:composeseeds-surjective}\lean{Obsmetrologia.PredicationCell.composeSeeds_surjective}\leanok\uses{thm:composeseeds, thm:predication}
Every predication cell \(\varphi\), on- or off-diagonal, equals
\(\texttt{composeSeeds}\,u\,v\) for some pair of utterances \(u, v\) ---
witnessed by \(u = (\varphi.\mathrm{subject}, \varphi.\mathrm{predicate})\)
and \(v = (\varphi.\mathrm{object}, \mathtt{default})\); the witness's second
claim is irrelevant, since \texttt{composeSeeds} never reads \(v\)'s claim.
This is a surjectivity statement over \emph{all} utterances, not one
restricted to a fixed recursion's own siblings --- that restriction is
\cref{thm:cellsof} below.
\end{theorem}

\begin{definition}[The cells composable from a recursion's own siblings]\label{thm:cellsof}\lean{Obsmetrologia.PredicationCell.cellsOf}\leanok\uses{thm:articulating-utterances, thm:composeseeds}
\(\texttt{cellsOf}\,a\) is the list of \(\texttt{composeSeeds}\,u\,v\) for
\(u, v\) both ranging over \(a.\mathrm{utterances}\) --- the provenance
restriction \cref{thm:composeseeds-surjective} does not itself supply: every
member of \(\texttt{cellsOf}\,a\) is composed from two of \(a\)'s own
siblings under \cref{rem:sibling-hypothesis}'s reading, not from an arbitrary
utterance drawn from outside.
\end{definition}

\begin{proposition}[The base case]\label{thm:cellsof-utter}\lean{Obsmetrologia.PredicationCell.cellsOf_utter}\leanok\uses{thm:cellsof}
A bare \(\texttt{utter}\,u\) has exactly one sibling, itself, so its only
composable cell is the diagonal seed: \(\texttt{cellsOf}\,(\texttt{utter}\,u)
= [\texttt{seed}\,u]\).
\end{proposition}

\begin{theorem}[The recursion's own diagonal seed is always drawn]\label{thm:seed-head-mem-cellsof}\lean{Obsmetrologia.PredicationCell.seed_head_mem_cellsOf}\leanok\uses{thm:cellsof, thm:head-mem-utterances, thm:composeseeds-self}
\(\texttt{seed}\,a.\mathrm{head} \in \texttt{cellsOf}\,a\) for every
recursion \(a\): the recursion always composes at least its own reflexive
seed from its own siblings, never needing anything from outside itself.
\end{theorem}

This section does not decide whether ``sibling'' should instead mean
something structural, only that the linear reading is formalized end to end.
Nor does \cref{thm:composeseeds-surjective} claim a \emph{unique} witnessing
pair for a given cell --- it produces one witness, not a bijection between
predication cells and pairs of utterances.

\begin{theorem}[The oriented cell articulating generates]\label{thm:oriented-cell}\lean{Obsmetrologia.Articulating.orientedCell, Obsmetrologia.Articulating.orientedCell_utter, Obsmetrologia.Articulating.orientedCell_subject, Obsmetrologia.Articulating.orientedCell_predicate, Obsmetrologia.Articulating.orientation_generated, Obsmetrologia.Articulating.converse_orientedCell_pair, Obsmetrologia.Articulating.subjProj_orientedCell, Obsmetrologia.Articulating.objProj_orientedCell_pair, Obsmetrologia.Articulating.orientedCell_pair_eq_composeSeeds, Obsmetrologia.Articulating.orientedCell_mem_cellsOf, Obsmetrologia.Articulating.tokenEdge}\leanok\uses{thm:composeseeds, thm:head-mem-utterances, thm:predication-seed, thm:converse-orientation}
The spine's predicating morphism reads only the head of an articulating, so every
cell it generates is the reflexive seed; \cref{thm:composeseeds} supplies the
off-diagonal cells as compositions of any two seeds from the recursion's flat pool
of utterances, under the linear-sibling hypothesis, and leaves open which
composition the recursion itself determines. The recursion \texttt{pair u a} is the
directed asymmetry: the head articulated over the tail. Its \emph{oriented cell} is
$(u.\mathrm{token},\,u.\mathrm{claim},\,a.\mathrm{head}.\mathrm{token})$: the seed
of the head composed with the seed of the utterance immediately below it, the one
composition the recursion picks out without a choice from the pool
(\texttt{orientedCell\_pair\_eq\_composeSeeds}, and it is a member of
\texttt{cellsOf}). On the floor it is the reflexive seed; at depth one orientation
is generated (\texttt{orientation\_generated}); the converse of a generated cell is
generated, by the articulation with the tokens exchanged and the claim carried; the
two projections recover the reflexive seeds of the two ends; and the token edge reads
only the head and the tail's head. The cells articulating generates are thus the
pairs of tokens with a claim between them, the shape a causal set has, reached from
articulating rather than taken as a given DAG. The
oriented cell is what \emph{measuring} reads --- \texttt{measure k trace} collects
witnesses by each articulating's own \texttt{orientedCell} --- so it is
articulating's generated cells, oriented, that a trace is measured against. Every
declaration is axiom-free.
\end{theorem}
\begin{theorem}[The spine predicates the oriented cell]\label{thm:spine-predicate-oriented}\lean{Obsmetrologia.Spine.predicate, Obsmetrologia.Spine.predicate_eq_orientedCell, Obsmetrologia.Spine.predicate_utter, Obsmetrologia.Spine.predicate_oriented}\leanok\uses{thm:oriented-cell}
The spine's predicating morphism is the oriented cell, the head articulated over
the tail, so the literal composite and the measured composite predicate alike. On
the floor it is the reflexive seed of the one utterance, so nothing that articulated
a single utterance reads differently; at depth one its subject and object can differ.
\end{theorem}
